
The combination of $m_W$ measurements is performed by first replicating each experiment's 
combination of fit results within the experiment, applying any relevant $\delta m_W$ shifts, 
and then combining across experiments.

The CDF individual $m_W$ values using the \mT, \pTl, and \met~distributions in the electron 
and muon channels are combined using the reported uncertainties and correlations, giving the
results shown in Table~\ref{tbl:cdfval}.  The CDF measurement applies $\delta m_W$ values of
$3.3$, $3.6$, and $3.0$~MeV, respectively, to fits to {\textsc{ResBos-C}} with the CTEQ6M PDF
set.  We remove these $\delta m_W$ corrections and add 
$\delta m_W^{\scriptsize\textrm{pol}} + \delta m_W^{\scriptsize\textrm{gen}} = -11.1$,
$-11.8$, and $-15.7$~MeV to the \mT, \pTl, and \met~results, respectively,
corresponding to the {\textsc{ResBos2}} calculation of leptonic angular distributions described
in Sec.~\ref{sec:polarization} and the removal of the $\ell\nu$ invariant mass requirement
discussed in Sec.~\ref{sec:resonance}.  Finally, a shift to the target PDF set is applied.
For the NNPDF3.1 PDF set this procedure
gives a combined CDF value of $m_W = 80432.1 \pm 9.4$~MeV, which is consistent with the published
CDF value of $m_W = 80433.5 \pm 9.4$~MeV within the uncertainty of the procedure.  

The principal component analysis used by CDF to reduce statistical effects in the PDF uncertainty
evaluation is not used in the combination with other experiments, since different measurements 
would give different principal components and complicate the correlation evaluations for the
combination.  Instead the Hessian sets provided by the NNPDF collaboration are used to estimate
the PDF uncertainty.  The combined CDF value with this uncertainty is labelled
``Combined ($\sigma_{\scriptsize\textrm{PDF}} = 6.6$~MeV)'' in Table~\ref{tbl:cdfval} and
corresponds to the entry labelled "NNPDF 3.1" in Table~\ref{tbl:tevatron_opt2} in
Sec.~\ref{sec:results}.

\begin{table*}[tbp]
  \centering
\begin{tabular*}{\linewidth}{l @{\extracolsep{\fill}}cccccc}
\hline \hline
                & Published $m_W$ & Input $m_W$ & $\delta m_W^{\scriptsize\textrm{pol}}$ & $\delta m_W^{\scriptsize\textrm{gen}}$
                & $\delta m_W^{\scriptsize\textrm{PDF}}$ & LHC-TeV MWWG $m_W$ 
                \\
\hline
\mT($e,\nu$)    & 80429.1     & 80425.8     & $-9.5$     & $-1.6$       &  $13.1$          & 80427.8 \\
\pTl($e$)       & 80411.4     & 80407.8     & $-8.4$     & $-3.4$       &  $14.0$          & 80410.0 \\
\met($e$)       & 80426.3     & 80423.3     & $-12.5$    & $-3.2$       &  $15.1$          & 80422.7 \\
\mT($\mu,\nu$)  & 80446.1     & 80442.8     & $-9.5$     & $-1.6$       &  $13.1$          & 80444.8 \\
\pTl($\mu$)     & 80428.2     & 80424.6     & $-8.4$     & $-3.4$       &  $14.0$          & 80426.8 \\
\met($\mu$)     & 80428.9     & 80425.9     & $-12.5$    & $-3.2$       &  $15.1$          & 80425.3 \\
\hline
Combined ($\sigma_{\scriptsize\textrm{PDF}} = 3.9$~MeV)  & 80433.5  & &   &                &          & 80432.1 \\
Combined ($\sigma_{\scriptsize\textrm{PDF}} = 6.6$~MeV)  &          & &   &                &          & 80433.3 \\
\hline \hline
\end{tabular*}
 \caption{Published CDF values and input values to the combination, where the latter correspond to the results 
 obtained from the direct CDF fits to {\textsc{ResBos-C}} with the CTEQ6M PDF set.  The combination procedure applies
 shifts to these results to update to the {\textsc{ResBos2}} calculation
 ($\delta m_W^{\scriptsize\textrm{pol}}$ and $\delta m_W^{\scriptsize\textrm{gen}}$)
 and a shift to update to the NNPDF3.1 PDF set ($\delta m_W^{\scriptsize\textrm{PDF}}$).  The total statistical and
 systematic uncertainties on the shifts are 1.2, 1.1, and 2.1 MeV for the \mT, \pTl, and \met~fits respectively.
 The combined value
 is consistent with that obtained by CDF when using the PDF uncertainties determined by CDF, labelled
 ``Combined ($\sigma_{\scriptsize\textrm{PDF}} = 3.9$~MeV)''.  When combining the result with other measurements, the
 uncertainty is evaluated using NNPDF3.1 eigenvectors to give the result labelled
 ``Combined ($\sigma_{\scriptsize\textrm{PDF}} = 6.6$~MeV)''.  The difference is due to a change in the weight of
 each fit distribution.  All units are in MeV.
}
\label{tbl:cdfval}
\end{table*}

The individual D0 measurements with the \mT~and \pTl~distributions using data sets corresponding 
to 1.1~fb$^{-1}$ and 4.3~fb$^{-1}$ are combined using the reported uncertainties to give the result 
$m_W = 80375.1 \pm 23.1$~MeV, which is the value quoted by D0 rounded to the nearest MeV.  Before 
combining with other measurements, a number of shifts are applied.  First, a shift from CTEQ6.1 to
CTEQ6.6 is applied to the measurement based on 1.1~fb$^{-1}$ of integrated luminosity.  Shifts of 
$\delta m_W^{\scriptsize\textrm{pol}} = -6.4$, $-6.9$, and $-15.8$ are applied to the \mT, \pTl,
and \met fit results, respectively, to update the {\textsc{ResBos-CP}} leptonic angular distributions to
those of {\textsc{ResBos2}}.  Finally, a $\delta m_W^{\mathrm\Gamma}$ shift adjusts ${\mathrm\Gamma_W}$
to that of the SM prediction.  The result with these shifts and the published D0 PDF uncertainty
of $\approx 11$~MeV is labelled ``Combined ($\sigma_{\scriptsize\textrm{PDF}} = 11$~MeV)'' in
Table~\ref{tbl:d0val}.  The value with uncertainties updated to those calculated with {\textsc{Wj-MiNNLO}}
and CTEQ6.6 is $m_W = 80377.9 \pm 25.5$~MeV and is labelled ``Combined ($\sigma_{\scriptsize\textrm{PDF}} = 15.1$~MeV)''
in the table.

\begin{table*}[tbp]
  \centering
\begin{tabular*}{\linewidth}{l @{\extracolsep{\fill}}ccccc}
\hline \hline 
  & Published $m_W$   & $\delta m_W^{\scriptsize\textrm{pol}}$ & $\delta m_W^{\scriptsize\textrm{PDF}}$
  & $\delta m_W^{\mathrm\Gamma}$ & LHC-TeV MWWG $m_W$ \\
\hline
Run 2a \mT($e,\nu$)  & 80401          & $-6.4$           & $14.3$         &  1.4 & 80410.3      \\
Run 2a \pTl($e$)     & 80400          & $-6.9$           & $16.7$         &  1.4 & 80411.2      \\
Run 2a \met($e$)     & 80402          & $-15.8$          & $17.5$         &  1.4 & 80405.1      \\
Run 2b \mT($e,\nu$)  & 80371          & $-6.4$           &                &  1.5 & 80366.1      \\
Run 2b \pTl($e$)     & 80343          & $-6.9$           &                &  1.5 & 80337.6     \\
\hline
Combined ($\sigma_{\scriptsize\textrm{PDF}} = 11$~MeV)           & 80375      &         & &      & 80373.4      \\ 
Combined ($\sigma_{\scriptsize\textrm{PDF}} = 15.1$~MeV)         &            &         & &      & 80377.9      \\ 
\hline \hline
\end{tabular*}
  \caption{Published D0 values corresponding to the CTEQ6M (Run 2a) and CTEQ6.6 (Run 2b) PDF sets, 
along with the following shifts: modifying the leptonic angular distributions to match those of
{\textsc{ResBos2}} ($\delta m_W^{\scriptsize\textrm{pol}}$);
modifying the Run 2a result to correspond to the CTEQ6.6 PDF set ($\delta m_W^{\scriptsize\textrm{PDF}}$); 
and modifying the $W$ boson width to the Standard Model prediction ($\delta m_W^{\mathrm\Gamma}$).
 The total statistical and systematic uncertainties on the shifts are 1.2, 1.2, and 2.3 MeV for the
 \mT, \pTl, and \met~fits respectively.
The combined result with the published D0 PDF uncertainty obtained using {\textsc Pythia} and the CTEQ6.1
PDF set is labelled ``Combined ($\sigma_{\scriptsize\textrm{PDF}} = 11$~MeV)'', and the result with
PDF uncertainties updated to those of CTEQ6.6 calculated with {\textsc{Wj-MiNNLO}} is labelled
``Combined ($\sigma_{\scriptsize\textrm{PDF}} = 15.1$~MeV)''.  The results differ due to different
weights of the individual fits to kinematic distributions.  All units are in MeV.  }
\label{tbl:d0val}
\end{table*}

The ATLAS measurement is reproduced using the parameterized simulation to give a value of 
$m_W = 80369.7 \pm 18.5$~MeV, which is within a few tenths of an MeV of the published result.  
A $\delta m_W^{\mathrm\Gamma} = 0.7$~MeV correction is added to update ${\mathrm\Gamma}_W$, and further
$\delta m_W^{\scriptsize\textrm{PDF}}$ shifts are applied to provide the central value for the
target PDF set.  All LHCb $m_W$ values for 
the combination are determined from a direct fit to the data using the target PDF set, so only a 
$\delta m_W^{\mathrm\Gamma} = -0.7$~MeV shift is applied to update the value of the $W$ boson width.
